\section{Projection into commercial coordinates}\label{projection-into-commercial-coordinates}

\subsection{1. Why a reduced model is necessary}\label{why-a-reduced-model-is-necessary}

The structural state contains distributions of trees, logs, processors, and
operators, together with a posterior over unknown parameters. Current regional
data identify only a subset of these objects.

Price, scale, cost, and break-even analysis uses a finite-dimensional
\textbf{projection} of the structural state.

The reduced model should be read as

\[
Z_t=\Pi_{\rm red}(\mu_t,\nu_t,\kappa_t,\pi_t),
\]

where \texttt{Pi\_red} keeps only the coordinates needed for today's commercial
questions.

\subsection{2. Asset-side value envelope}\label{asset-side-value-envelope}

For asset account \texttt{i}, the paper defines a value envelope

\[
V_{A,i}=A_i^{ha}\left(S_i\delta_i^{op}
+V_i^{asset}\lambda_i^{loss}\right).
\]

Interpretation:

\begin{itemize}
\tightlist
\item
  \texttt{A\_i\^{}ha} is managed area;
\item
  \texttt{S\_i} is annual forestry spend per hectare;
\item
  \texttt{delta\_i\^{}op} is an information-enabled proportional operating improvement;
\item
  \texttt{V\_i\^{}asset} is standing asset value per hectare;
\item
  \texttt{lambda\_i\^{}loss} is avoidable annual loss/risk.
\end{itemize}

Information creates value through operational efficiency and avoided loss. The
revenue model applies a rate to managed spend and captures part of this envelope.

\subsection{3. Asset monetisation and adoption}\label{asset-monetisation-and-adoption}

Let \texttt{C\_i\^{}A} be an adoption indicator and define the adopted managed-spend base

\[
B_A=\sum_i C_i^A A_i^{ha}S_i.
\]

Asset revenue is

\[
R_A=m_A B_A.
\]

For one homogeneous segment, the paper uses logistic adoption

\[
a(p)=\frac{1}{1+\exp(\eta_Ap-\alpha_A)}.
\]

If the monetisable base is \texttt{B}, revenue is

\[
R(p)=Bp\,a(p).
\]

The log revenue is

\[
\log R(p)=\log B+\log p-\log(1+e^{\eta_Ap-\alpha_A}).
\]

Differentiating twice gives

\[
\frac{d^2}{dp^2}\log R(p)
=-\frac1{p^2}-\eta_A^2a(p)(1-a(p))<0.
\]

Therefore revenue is strictly log-concave and has a unique maximizer for
positive price.

The first-order condition is

\[
\eta_Ap(1-a(p))=1.
\]

Solving yields

\[
p^*=\frac{1+W(e^{\alpha_A-1})}{\eta_A},
\]

where \texttt{W} is the principal Lambert W branch.

This result is local to the simplified one-segment model. Mixtures of customer
segments, fixed packages, network effects, or dynamic learning can destroy the
one-dimensional global shape.

\subsection{4. Processor-side monetisation}\label{processor-side-monetisation}

The structural matching layer outputs matched tonnes \texttt{M\_j}. If processor \texttt{j}
adopts the service and pays \texttt{f\_j} per matched/visible tonne,

\[
R_S=\sum_j C_j^S f_j M_j.
\]

\texttt{M\_j} is the quantity remaining after compatibility, spatial cost, and capacity
constraints are applied to regional wood supply.

Today's simpler ``commercially visible supply'' assumption is therefore an
approximation to an operator that has a precise theoretical meaning.

\subsection{5. Cost structure}\label{cost-structure}

The paper uses

\[
C(H,N_A,N_P,n_V,W_D,G_{act})
=F+\sum_{g\in G_{act}}F_g+a_HH^{\beta_H}
+b_AN_A+b_PN_P+c_Vn_V+c_DW_D+\epsilon_C.
\]

Components include:

\begin{itemize}
\tightlist
\item
  common fixed cost \texttt{F};
\item
  geography-specific fixed blocks \texttt{F\_g};
\item
  variable coverage cost \texttt{a\_H\ H\^{}\{beta\_H\}};
\item
  account-level service costs;
\item
  verification cost;
\item
  data/compute workload;
\item
  residual cost shock.
\end{itemize}

This form is intentionally able to separate \textbf{fixed-cost dilution} from genuine
sublinear marginal scaling.

\subsection{6. Fixed-cost dilution}\label{fixed-cost-dilution}

If

\[
C(H)=F+aH^\beta,
\]

then local total-cost elasticity is

\[
\frac{d\log C}{d\log H}
=\beta\frac{aH^\beta}{F+aH^\beta}<\beta.
\]

So even if the variable component is nearly linear (\texttt{beta} near one), total
cost can appear strongly sublinear at small-to-medium scale because the fixed
cost is being spread over more hectares.

Three declining average-cost observations leave the variable exponent
unidentified. Additional scale observations or stronger cost restrictions are
required.

\subsection{7. Operating leverage}\label{operating-leverage}

Let total recurring revenue be \texttt{R\_H} and expected recurring cost \texttt{C\_H}. Define
local elasticities

\[
\beta_R(H)=\frac{d\log E[R_H]}{d\log H},\qquad
\beta_C(H)=\frac{d\log E[C_H]}{d\log H}.
\]

Positive operating leverage requires

\[
\beta_R(H)>\beta_C(H).
\]

Positive operating leverage requires revenue to grow faster than recurring cost
over the relevant range.

\subsection{8. Break-even frontier}\label{break-even-frontier}

At fixed scale,

\[
m_AB_A+f_SB_S=C,
\]

where \texttt{B\_S} is the adopted matched-tonnage base.

For fixed bases and nonnegative rates, the set

\[
\{(m_A,f_S):m_AB_A+f_SB_S\ge C\}
\]

is a convex half-space. This makes simple trade-off analysis possible: stronger
asset monetisation can compensate for weaker processor monetisation and vice
versa.

The convexity disappears if the bases themselves respond nonlinearly to price,
network effects, or participation thresholds.

\subsection{9. A risk multiplier}\label{a-risk-multiplier}

Let \texttt{R\_H\textgreater{}0} and \texttt{C\_H\textgreater{}0} be random recurring revenue and cost. For target
probability \texttt{tau}, define the revenue multiplier

\[
M_H^\tau=\inf\{k>0:P(kR_H-C_H\ge0)\ge\tau\}.
\]

Because revenue is positive,

\[
M_H^\tau=Q_\tau(C_H/R_H),
\]

where \texttt{Q\_tau} is the left quantile.

This gives an interpretable stress metric: by what multiplicative factor would
the current revenue distribution need to improve to achieve a desired chance of
nonnegative recurring surplus?

\subsection{10. Scope of the projection}\label{scope-of-the-projection}

The reduced model provides conditional projections under stated adoption,
matched-supply, and cost assumptions. It supplies no causal effect of hectares on
profitability. New observations replace coarse inputs with structural estimates.
